\documentclass[10pt,a4paper]{amsart}
\usepackage[margin=19mm]{geometry}
\usepackage{amsmath,amssymb,mathtools}
\usepackage[scr=rsfs,cal=euler]{mathalfa}
\usepackage{xfrac}
\usepackage[unicode,hidelinks,bookmarks=false]{hyperref}
\newcommand{\ie}{i.e.\ }
\newcommand{\cf}{cf.\ }
\newcommand{\resp}{respectively\ }
\newcommand{\Subsection}{\subsection}
\providecommand{\nobreakdash}{\nobreak\hbox{-}}
\setlength{\emergencystretch}{2em}
\multlinegap0pt
\usepackage[shortlabels]{enumitem}
\usepackage{braket}
\usepackage{mathtools}
\usepackage{amssymb}
\usepackage{csquotes}
\usepackage{tikz}
\newtheorem{proposition}{Proposition}
\newtheorem{corollary}{Corollary}
\newcommand{\N}{\mathbb{N}}
\title{A Centroid Framework for Operator-Valued Haagerup Inequalities}
\author{Patrick Oliveira Santos}
\date{Source: arXiv:2608.26643v1 (CC BY 4.0)}
\begin{document}
\maketitle

\noindent\emph{Source and licence.} A Centroid Framework for Operator-Valued Haagerup Inequalities. Version arXiv:2608.26643v1 (CC BY 4.0), distributed by arXiv under the Creative Commons Attribution 4.0 International licence, http://creativecommons.org/licenses/by/4.0/, as recorded in the arXiv licence metadata for this exact version and shown on the abstract page https://arxiv.org/abs/2608.26643v1. Full licence notice: \texttt{licenses/arxiv-2608.26643-CC-BY-4.0.txt}.\par\smallskip

\noindent\textit{Editorial scope.} Counted unit: the article's proposition approximation (source0.tex lines 1069--1079). The article's complete original proof of it is delivered with it (source0.tex lines 1209--1269, one \texttt{proof} environment of 61 lines, which the article places away from the statement). No mathematics is written, repaired or re-derived: the statement and the proof are the article's own lines, in the article's own notation, and no retained line is substituted or re-flowed.  Retained uncounted as the source-local context the proof needs: lines 1160--1170.  Nothing was omitted from any retained span, so the package carries no omission record.  No retained line contains a citation, so the wrapper carries no bibliography and no bibliography entry.  No retained line contains a figure environment, an image-inclusion command or drawing code, so no image is delivered and the package declares no asset dependency.  Editorial formatting only, in addition: an AMS \texttt{amsart} preamble that restates the article's own declarations and macros used by the retained lines, namely the environments \texttt{proposition}, \texttt{corollary} on separate counters, together with the standard packages those lines need.  The retained environments print in this document as Proposition 1, Corollary 1 in order of appearance, whereas the article prints the same items with the numbering of its own full document.  The complete native source of the article as distributed in the arXiv e-print for this exact version is bundled unchanged as \texttt{sources/208653/source0.tex}.
\begin{corollary}\label{corollary: iterative procedure}
    There exist constants $C_1,C_2,C_3>0$, depending only on $T$, the coarse parameters, and the rank, with the following property. Let $Z\subseteq X$, $a\in X$, and $E\ge 0$ satisfy
    \begin{align*}
        a\in [x_0,z]_E\qquad\text{for every }z\in Z.
    \end{align*}
    Then there exists $h\in X$ such that
    \begin{align*}
        &h\in [x_0,z]_{\theta E+C_1}\qquad\text{for every }z\in Z,\\
        &d(a,h)\le C_2E+C_3.
    \end{align*}
\end{corollary}

\begin{proposition}\label{proposition: coarse approximation}
    Let $(X,d,\mu)$ be a coarse median metric space of rank at most $\nu$, with base point $x_0$. Then there are constants $E^*,K_1,K_2$ depending only on $\nu$ and the coarse parameters such that, for every $Z\subseteq X$ and $a\in X$ satisfying
    \begin{align*}
        R(a,Z):=\sup_{z\in Z}d(a,z)<\infty,
    \end{align*}
    there exists $h\in X$ with the properties
    \begin{enumerate}[label=(P\arabic*)]
        \item\label{property: inclusion with E*} $h\in [x_0,z]_{E^*}$ for every $z\in Z$.
        \item\label{property: distance bound with K_1K_2} $d(a,h)\le K_1R(a,Z)+K_2$.
    \end{enumerate}
\end{proposition}

\begin{proof}[Proof of Proposition~\ref{proposition: coarse approximation}]
    For every $z\in Z$,
    \begin{align*}
        d(a,\braket{x_0,z,a})=d(\braket{x_0,a,a},\braket{x_0,z,a})\le\rho R(a,Z)+h(0)=:E_0.
    \end{align*}
    Thus $a\in[x_0,z]_{E_0}$ for every $z\in Z$.
    Apply Corollary~\ref{corollary: iterative procedure} with $a_0=a$ and $E=E_0$. This yields $a_1\in X$ such that
    \begin{align*}
        &a_1\in [x_0,z]_{\theta E_0+C_1}\qquad\text{for every }z\in Z,\\
        &d(a_1,a)\le C_2E_0+C_3.
    \end{align*}
    Recursively, let $a_{n+1}\in X$ be the output of Corollary~\ref{corollary: iterative procedure} applied with $a=a_n$ and $E=E_n$. Then
    \begin{equation}
    \label{equation: recursive bounds}
    \begin{aligned}
        &a_{n+1}\in [x_0,z]_{E_{n+1}}\qquad\text{for every }z\in Z,\\
        &E_{n+1}=\theta E_n+C_1\\
        &d(a_{n+1},a_{n})\le C_2 E_n+C_3.
    \end{aligned}      
    \end{equation}
    Define
    \begin{align*}
        E^*=\max \left\{1,\frac{2C_1}{1-\theta}\right\};\quad \delta=\frac{1+\theta}{2}<1.
    \end{align*}
    Note that $E^*$ and $\delta$ depend only on $\nu$ and the coarse parameters.
    Whenever $E_n>E^*$,
    \begin{align*}
        E_n>\frac{2C_1}{1-\theta},
    \end{align*}
    and therefore $C_1< \frac{1-\theta}{2}E_n$. In particular,
    \begin{align*}
        E_{n+1}=\theta E_n+C_1< \theta E_n+\frac{1-\theta}{2}E_n=\delta E_n.
    \end{align*}
    Since $\delta<1$, define
    \begin{align*}
        N=\min\{n\in \N: E_n \le E^*\}<\infty.
    \end{align*}
    Then, for every $k<N$, $E_{k+1}\le \delta E_k$.
    Set $h=a_N$.
    By \eqref{equation: recursive bounds}, we have
    \begin{align*}
        h\in [x_0,z]_{E_N}\subseteq [x_0,z]_{E^*}\qquad\text{for every }z\in Z.
    \end{align*}
    This proves Property~\ref{property: inclusion with E*}. The triangle inequality gives
    \begin{align*}
        d(a,h)\le \sum_{k=0}^{N-1} d(a_k,a_{k+1}).
    \end{align*}
    For every $0\le k<N$, $E_k>E^*$. Hence, \eqref{equation: recursive bounds} yields
    \begin{align*}
        d(a_k,a_{k+1})\le C_2 E_k+C_3=\left(C_2+\frac{C_3}{E_k}\right)E_k\le \left(C_2+\frac{C_3}{E^*}\right)E_k.
    \end{align*}
    Consequently,
    \begin{align*}
        d(a,h)\le \left(C_2+\frac{C_3}{E^*}\right) \sum_{k=0}^{N-1}E_k.
    \end{align*}
    Since $E_{k+1}\le\delta E_k$, the geometric-series estimate gives
    \begin{align*}
        d(a,h)\le \left(C_2+\frac{C_3}{E^*}\right)E_0 \sum_{k=0}^{\infty}\delta^k=\frac{1}{1-\delta}\left(C_2+\frac{C_3}{E^*}\right)E_0.
    \end{align*}
    The constants $\delta,C_2,C_3,E^*$ depend only on the coarse parameters and $\nu$, while $E_0=\rho R(a,Z)+h(0)$. This proves Property~\ref{property: distance bound with K_1K_2}.
\end{proof}

\end{document}
